% Part: incompleteness
% Chapter: incompleteness-provability
% Section: fixed-point-lemma

\documentclass[../../../include/open-logic-section]{subfiles}

\begin{document}

\olfileid{inc}{inp}{fix}

\olsection{નિશ્ચિત-બિંદુ લેમ્મા}

\begin{explain}
નિશ્ચિત-બિંદુ લેમ્મા કહે છે કે
કોઈપણ !!{formula}~$!B(x)$ માટે
એવું !!a{sentence}~$!A$ છે કે
$\Th{T} \Proves !A \liff !B(\gn{!A})$,
જો $\Th{T}$ એ~$\Th{Q}$નું
વિસ્તરણ હોય. જૂઠા વાક્યના
કિસ્સામાં, આપણે $!A$ એ
~``$\gn{!A}$ અસત્ય છે''
એની સાથે (~$\Th{T}$માં
સાબિતિક્ષમ રીતે) સમકક્ષ
હોય એવું ઇચ્છીએ; એટલે
$\Gn{!A}$ એ અસત્ય
!!{sentence}ની ગોડેલ સંખ્યા
છે એવું વિધાન.
સાબિતીનો વિચાર સમજવા,
$!A$ માટે ક્વાઇનનું
અનૌપચારિક રૂપ સરખાવવું
ઉપયોગી છે: ``{}`પોતાનું
અવતરણ પહેલાં મૂકવાથી
અસત્ય મળે છે' પોતાનું
અવતરણ પહેલાં મૂકવાથી
અસત્ય મળે છે.'' કોઈ
અભિવ્યક્તિની પહેલાં તેનું
પોતાનું અવતરણ મૂકીને
વાક્ય રચવાની ક્રિયાને તે
અભિવ્યક્તિનું
\emph{વિકર્ણીકરણ} કહી શકાય;
મળેલું વાક્ય તેનું
વિકર્ણિત રૂપ છે.
એટલે `પોતાનું અવતરણ
પહેલાં મૂકવાથી અસત્ય મળે
છે'નું વિકર્ણિત રૂપ
``{}`પોતાનું અવતરણ પહેલાં
મૂકવાથી અસત્ય મળે છે'
પોતાનું અવતરણ પહેલાં
મૂકવાથી અસત્ય મળે છે'' છે.
હવે નોંધો કે ક્વાઇનનું
જૂઠું વાક્ય `અસત્ય મળે છે'નું
વિકર્ણિત રૂપ નથી, પણ
`પોતાનું અવતરણ પહેલાં
મૂકવાથી અસત્ય મળે છે'નું
છે. તેથી જૂઠું વાક્ય
મેળવવા જે ગુણધર્મનું
વિકર્ણીકરણ થાય છે તેમાં
જ વિકર્ણીકરણ સમાયેલું છે!{}

અંકગણિતની ભાષામાં એક
મુક્ત ચલવાળા !!a{formula}ની
ગોડેલ સંખ્યા ગણીએ અને
પછી તે ચલને એ ગોડેલ
સંખ્યાના પ્રમાણભૂત અંકપદથી
બદલીએ; આમ તેના અવતરણ
રચાય છે.
~$!E(x)$નું વિકર્ણિત
રૂપ $!E(\num{n})$ છે,
જ્યાં $n = \Gn{!E(x)}$.
(હવેથી $\num{\Gn{!E(x)}}$ને
$\gn{!E(x)}$ લખીશું.)
તેથી $!B(x)$ જો
``અસત્ય છે'' હોય,
તો ``પોતાનું અવતરણ પહેલાં
મૂકતાં અસત્ય મળે છે''
એનો અર્થ ``પોતાના
વિકર્ણિત રૂપની ગોડેલ
સંખ્યા પર લાગુ કરતાં
અસત્ય મળે છે'' થાય.
ધારો કે !!{formula}ની
ગોડેલ સંખ્યા~$n$ પરથી
તેના વિકર્ણિત રૂપની
ગોડેલ સંખ્યા ગણતા વિધેય
$\fn{diag}(n)$ માટે
આપણી પાસે
~$\Obj{diag}$ ચિહ્ન હોય.
તો $!E(x)$ને
$!B(\Obj{diag}(x))$
લખી શકીએ. અને ક્વાઇનના
જૂઠા વાક્યનું રૂપ તેનું
વિકર્ણીકરણ, એટલે
$!E(\gn{!E(x)})$ અથવા
$!B(\Obj{diag}(\gn{!B(\Obj{diag}(x))}))$,
થાય. અલબત્ત, હવે
$!B(x)$ બીજો કોઈપણ
ગુણધર્મ હોઈ શકે અને એ
જ રચના કામ કરશે.
અપૂર્ણતા પ્રમેય માટે
$!B(x)$ તરીકે
``$x$ એ~$\Th{T}$માં
!!{derivable} નથી'' લઈશું.
ત્યારે $!E(x)$નો અર્થ
``પોતાના વિકર્ણિત રૂપની
ગોડેલ સંખ્યા પર લાગુ
કરતાં $\Th{T}$માં
!!{derivable} ન હોય એવું
!!a{sentence} મળે છે''
થાય.

આને~$\Th{T}$માં ઔપચારિક
બનાવવા $\fn{diag}$ને
ઔપચારિક કરવાની રીત
જોઈએ. વિધેય
$\fn{diag}(n)$ સંગણનીય,
હકીકતમાં આદિમ પુનરાવર્તી
છે: $n$ જો
સૂત્ર~$!E(x)$ની ગોડેલ
સંખ્યા હોય, તો
$\fn{diag}(n)$ એ
~$!E(\gn{!E(x)})$ની
ગોડેલ સંખ્યા આપે છે.
(યાદ રાખો,
$\gn{!E(x)}$ એ
~$!E(x)$ની ગોડેલ સંખ્યાનું
પ્રમાણભૂત અંકપદ છે,
એટલે $\num{\Gn{!E(x)}}$.)
જો $\Obj{diag}$ એ
$\fn{diag}$નું નિરૂપણ
કરતું $\Th{T}$માંનું
વિધેય-ચિહ્ન હોય,
તો $!A$ તરીકે સૂત્ર
$!B(\Obj{diag}(\gn{!B(\Obj{diag}(x))}))$
લઈ શકીએ. નોંધો કે
\begin{align*}
\fn{diag}(\Gn{!B(\Obj{diag}(x))}) & =
\Gn{!B(\Obj{diag}(\gn{!B(\Obj{diag}(x))}))} \\
& = \Gn{!A}.
\end{align*}
માનો કે $\Th{T}$ નીચેનું
!!{derive} કરી શકે:
\[
\Obj{diag}(\gn{!B(\Obj{diag}(x))}) = \gn{!A},
\]
તો તે
$!B(\Obj{diag}(\gn{!B(\Obj{diag}(x))}))
\liff !B(\gn{!A})$
પણ !!{derive} કરી શકે.
પરંતુ ડાબી બાજુ
વ્યાખ્યા મુજબ~$!A$ છે.

અલબત્ત, સામાન્ય રીતે
$\Obj{diag}$ એ
$\Th{T}$નું વિધેય-ચિહ્ન
હશે નહીં; ખાસ કરીને
~$\Th{Q}$નું તો નથી જ.
પરંતુ $\fn{diag}$ સંગણનીય
હોવાથી કોઈ સૂત્ર
$!D_{\fn{diag}}(x,y)$ તેને
~$\Th{Q}$માં
\emph{નિરૂપે} છે.
તેથી $!B(\Obj{diag}(x))$
લખવાને બદલે
$\lexists[y][(!D_{\fn{diag}}(x,y) \land !B(y))]$
લખી શકીએ. બાકી ઉપર
રેખાંકિત કરેલી સાબિતી
ચાલે છે; ખરેખર, તે
~$\Th{Q}$માં જ ચાલે છે.
\end{explain}

\begin{lem}
\ollabel{lem:fixed-point}
$!B(x)$ એ માત્ર એક
મુક્ત ચલ~$x$ ધરાવતું
કોઈપણ સૂત્ર હોય.
તો એવું વાક્ય~$!A$ છે કે
$\Th{Q} \Proves
!A \liff !B(\gn{!A})$.
\end{lem}


\begin{proof}
$!B(x)$ આપેલું હોય ત્યારે
$!E(x)$ એ સૂત્ર
$\lexists[y][(!D_{\fn{diag}}(x,y) \land !B(y))]$
લો અને $!A$ એ તેનું
વિકર્ણિત રૂપ, એટલે
સૂત્ર $!E(\gn{!E(x)})$ લો.

$!D_{\fn{diag}}$ એ
$\fn{diag}$નું નિરૂપણ કરે
છે અને
$\fn{diag}(\Gn{!E(x)}) = \Gn{!A}$
હોવાથી $\Th{Q}$ નીચેનું
!!{derive} કરી શકે છે:
\begin{align}
  & !D_{\fn{diag}}(\gn{!E(x)}, \gn{!A}) \ollabel{repdiag1} \\
  & \lforall[y][(!D_{\fn{diag}}(\gn{!E(x)},y) \lif
  \eq[y][\gn{!A}])]. \ollabel{repdiag2}
\end{align}
હવે બતાવીએ કે
$\Th{Q} \Proves !A \liff !B(\gn{!A})$.
માત્ર તર્ક અને
~$\Th{Q}$માં !!{derivable}
હકીકતો વાપરીને
અનૌપચારિક દલીલ કરીશું.

પ્રથમ $!A$, એટલે
$!E(\gn{!E(x)})$ માનો.
$!E(x)$ની વ્યાખ્યા પરથી
$!E(\gn{!E(x)})$ એટલે
\[
\lexists[y][(!D_{\fn{diag}}(\gn{!E(x)},y) \land !B(y))].
\]
આવો~$y$ લો.
$!D_{\fn{diag}}(\gn{!E(x)},y)$
હોવાથી \olref{repdiag2}
મુજબ $y = \gn{!A}$.
તેથી $!B(y)$ પરથી
~$!B(\gn{!A})$ મળે છે.

હવે $!B(\gn{!A})$ માનો.
\olref{repdiag1} મુજબ
\begin{align*}
& !D_{\fn{diag}}(\gn{!E(x)}, \gn{!A}) \land !B(\gn{!A}).
\intertext{આથી મળે છે કે}
& \lexists[y][(!D_{\fn{diag}}(\gn{!E(x)},y) \land !B(y))].
\end{align*}
પરંતુ એ જ
$!E(\gn{!E(x)})$,
એટલે~$!A$ છે.
\end{proof}

\begin{digress}
આ સાબિતીને સંગણનીયતા
સિદ્ધાંતના નિશ્ચિત-બિંદુ
લેમ્માની સાબિતી સાથે
સરખાવો. તફાવત એ છે કે
અહીં પોતાને આધારે
\emph{વિધાન} વ્યાખ્યાયિત
કરવું છે, જ્યારે ત્યાં
પોતાને આધારે
\emph{વિધેય} વ્યાખ્યાયિત
કરવું હતું. આ ભેદ સિવાય
વિચાર ખરેખર એક જ છે.
\end{digress}

\begin{prob}
  !!^a{formula}~$!A(x)$ને
  \emph{સત્યની વ્યાખ્યા}
  ત્યારે કહીએ જ્યારે બધાં
  !!{sentence}~$!B$ માટે
  $\Th{Q}
  \Proves !B \liff !A(\gn{!B})$
  હોય. નિશ્ચિત-બિંદુ લેમ્મા
  વાપરીને બતાવો કે કોઈ
  !!{formula} સત્યની
  વ્યાખ્યા નથી.
\end{prob}

\end{document}
